% ============================================================================
% Section 3.1 — Integers and Their Opposites (VA SOL: rational number emphasis)
% Modified topic for Virginia
% ============================================================================
\section{Integers and Their Opposites}

\begin{stepGoal}
\begin{itemize}
    \item Identify integers, opposites, and absolute values on a number line
    \item Compare and order rational numbers using inequality symbols
\end{itemize}
\end{stepGoal}

\begin{stepVocabBox}
    \vocabItem{Integer}{A whole number or its negative: $\ldots, -3, -2, -1, 0, 1, 2, 3, \ldots$}
    \vocabItem{Opposite}{The same distance from $0$ on the other side of the number line.}
    \vocabItem{Absolute Value}{The distance from $0$, written $|n|$. Always $\geq 0$.}
\end{stepVocabBox}

\begin{stepsCard}[How to Compare and Order Rational Numbers]
    \stepItem{Convert all numbers to the \textbf{same form} (decimals are usually easiest).}
    \stepItem{Place the numbers on a \textbf{number line} — left is less, right is greater.}
    \stepItem{Write the order using $<$ or $>$, or list from \textbf{least to greatest}.}
\end{stepsCard}

% --- Worked Example 1 ---
\begin{stepExample}{Find the opposite and absolute value of $-7$.}
    \stepShow{1} $-7$ as a decimal is already $-7.0$.

    \stepShow{2} $-7$ is $7$ units left of $0$. Its opposite ($7$) is $7$ units right.

    \stepShow{3} Opposite: $7$. Absolute value: $|-7| = 7$. Check: $-7 + 7 = 0$ $\checkmark$
    \stepResult{Opposite: $7$ \quad $|-7| = 7$}
\end{stepExample}

% --- Worked Example 2 ---
\begin{stepExample}{Order from least to greatest: $-1.5, \;\frac{3}{4}, \;-\frac{7}{4}, \;0.2, \;-1$}
    \stepShow{1} Convert to decimals: $-1.5, \;0.75, \;-1.75, \;0.2, \;-1$.

    \stepShow{2} On the number line, left to right: $-1.75, \;-1.5, \;-1, \;0.2, \;0.75$.

    \stepShow{3} $-\frac{7}{4} < -1.5 < -1 < 0.2 < \frac{3}{4}$.
    \stepResult{$-\dfrac{7}{4}, \;-1.5, \;-1, \;0.2, \;\dfrac{3}{4}$}
\end{stepExample}

\mascotStepTip{To compare fractions and decimals, convert them all to decimals first. It makes ordering much easier!}

% ============================================================================
% PRACTICE
% ============================================================================
\newpage
\begin{stepPractice}{Integers and Their Opposites Practice (VA)}
    \resetProblems

    \stepPracticeHeader[step1Color]{Convert to the Same Form}
    \begin{multicols}{2}
    \prob Write $-\frac{2}{3}$ as a decimal. \answerBlank[2cm]
    \answer{$\approx -0.667$}
    \prob Write $0.25$ as a fraction. \answerBlank[2cm]
    \answer{$\frac{1}{4}$}
    \end{multicols}

    \stepPracticeHeader[step2Color]{Compare Using $<$, $>$, or $=$}
    \begin{multicols}{2}
    \prob $-\frac{2}{3}$ \answerBlank[1.5cm] $-0.5$
    \answer{$<$}
    \prob $|-3|$ \answerBlank[1.5cm] $|2|$
    \answer{$>$}
    \end{multicols}

    \stepPracticeHeader[step3Color]{Order from Least to Greatest}
    \prob $\frac{1}{2}, \;-0.75, \;-\frac{1}{4}, \;0.3, \;-1$

    \answerBlank[6cm]
    \answer{$-1, \;-0.75, \;-\frac{1}{4}, \;0.3, \;\frac{1}{2}$}

    \wordProblem{A football team gains $6$ yards, then loses $6$ yards. Use integers to show the result. What is the opposite of $-6$?}{yards}
    \answer{$6 + (-6) = 0$ yards. The opposite of $-6$ is $6$.}
\end{stepPractice}
